---
title: "泛函分析作业250929"
date: 2025-09-29
collections: "泛函分析作业"
tags: ["泛函分析"]
categories: ["泛函分析"]
---

\problem[1.3.5] 设 $M$ 是 $C[a,b]$ 中的有界集, 求证: 集合 $$\{F(x)=\int_a^x f(t)\td t|f\in M\}$$ 是列紧集.

\begin{proof}
	根据微积分基本定理 $F(x)$ 也是 $[a,b]$ 上的连续函数.

	所以考虑使用 Arzela-Ascoli 定理. 即需要证明 $F$  一致有界且等度连续.

	"一致有界": 根据 $M$ 的有界性, 存在 $C>0$ 使得 $|f(x)|\leqslant C,\forall f\in M, x\in[a,b]$. 所以有 $$|F(x)|=|\int_a^x f(t)\td t|\leqslant\int_a^x|f(t)|\td t\leqslant (b-a)C$$ 所以 $F(x)$ 一致有界.

	"等度连续": 对任意 $\varepsilon>0$, 取 $\delta=\frac{\varepsilon}{C}$, 只要 $|x_1-x_2|<\delta$ 就有 $|F(x_1)-F(x_2)|<C|x_1-x_2|<\varepsilon,\forall F$. 从而也是等度连续.

	综上由 AA 定理可知该集合是列紧集.
\end{proof}

\problem[1.3.6] 设 $E=\{\sin nt\}_{n=1}^\infty$, 求证: $E$ 在 $C[0,\pi]$ 中不是列紧的.

\begin{proof}
	根据 AA 定理, 只需证 $E$ 不满足等度连续即可.

	$\forall \delta>0$, 存在 $k\in\N$ 使得 $\frac \pi{2^k}<\delta$, 从而对于 $x_1=0,x_2=\frac \pi {2^k}, f(t)=\sin 2^{k-1}t$, 有 $|x_1-x_2|<\delta$ 但 $|f(x_2)-f(x_1)|=\sin\frac\pi{2}=1$ 所以不可能等度连续. 即 $E$ 不是列紧的.
\end{proof}

\problem[1.3.7] 求证: $S$ 空间的子集 $A$ 列紧的充要条件是: $\forall n\in \N,\exists C_n>0$, 使得对 $\forall x=(\xi_1,\xi_2,\cdots,\xi_n,\cdots)\in A$, 有 $|\xi_n|\leqslant C_n(n=1,2,\cdot)$.

$S$ 空间: 在数列上定义距离 $\rho(x,y)=\sum\limits_{k=1}^\infty \frac{1}{2^k}\cdot\frac{|\xi_k-\eta_k|}{1+|\xi_k-\eta_k|}$.

\begin{proof}
	"$\Rightarrow$": 反设条件不成立, 即 $\exists n\in\N$, 不存在 $C_n>0$ 使得 $\xi_n$ 有界. 也就是说第 $n$ 个分量无界. 那么就可以取一列 $\{x^{(m)}\}$ 满足 $\xi_n^{(m)}>\xi_n^{(m-1)}+1$, 从而就有 $\rho(x^{(m)},x^{(m_2)})\geqslant \frac 1 {2^n}\cdot\frac{|\xi_n^{(m)}-\xi_n^{(m_2)}|}{1+|\xi_n^{(m)}-\xi_n^{(m_2)}|}\geqslant\frac{1}{2^{n+1}}$. 所以不存在收敛子列, 与 $A$ 列紧矛盾. 故该条件成立.

	"$\Leftarrow$": 由于 $S$ 空间完备, 故考虑证明在该条件下 $A$ 是完全有界集. 即证明存在有穷 $\varepsilon$-网.

	$\forall \varepsilon>0$, 取 $N=\lceil -\log_2 \varepsilon \rceil + 1$, 那么对于度量中 $\sum\limits_{k=N+1}^\infty \frac{1}{2^k}\cdot\frac{|\xi_k-\eta_k|}{1+|\xi_k-\eta_k|}<\sum\limits_{k=N+1}^\infty\frac{1}{2^k}<\frac{\varepsilon}{2}$.

	而对于前 $N$ 项, 定义集合 $D=\{(\xi_1,\xi_2,\ldots,\xi_N)|\exists x\in A, s.t. x_i=\xi_i\ i=1,2,\ldots,N\}$.

	根据条件每一分量都是有界的, 又只有有限维所以 $D$ 作为 $\mathbb K^N$ 的子集是有界的, 所以是完全有界的, 即存在有穷 $\frac{\varepsilon}{2}$-网 $N$. 而根据 $D$ 的定义. 这个网中的每个元素都对应 $A$ 中的一个或多个元素, 我们取其中任意一个作为该网中元素在 $A$ 中的代表元, 即存在有限集 $D'=\{(\xi_1,\xi_2,\ldots,\xi_N,\ldots)\}\subset A$ 满足 $\forall x\in A, \exists y\in D'$, 使得 $(\sum\limits_{i=1}^N |\xi_i-\eta_i|^2)^{1/2}<\varepsilon/2$, 从而有 $\forall i\in[1,N], |\xi_i-\eta_i|<\varepsilon/2$.

	所以 $\rho(x,y)=\sum\limits_{k=1}^N+\sum\limits_{k=N+1}^\infty<\frac\varepsilon 2+\frac\varepsilon 2=\varepsilon$.

	故 $D'$ 就是 $A$ 的有穷 $\varepsilon$-网. 从而 $A$ 是完全有界集, 进而是列紧集.
\end{proof}

\problem[1.3.8] 设 $(\mathscr X,\rho)$ 是度量空间, $M$ 是 $\mathscr X$ 中的列紧集, 映射 $f:\mathscr X\to M$ 满足 $$\rho(f(x_1),f(x_2))<\rho(x_1,x_2),\quad (\forall x_1,x_2\in \mathscr X,x_1\neq x_2)$$ 求证: $f$ 在 $\mathscr X$ 中存在唯一不动点.

\begin{proof}
	记 $d=\inf\{\rho(x,f(x))|x\in M\}$.
	
	由下确界的定义, 存在点列 $\{x_n\}$ 满足 $\rho(x_n,f(x_n))<d+\frac 1 n$. 由于 $M$ 是列紧集, 故存在收敛子列 $\{x_{n_k}\}$ 且满足 $\rho(x_{n_k},f(x_{n_k}))<d+\frac 1{n_k}$, 设其收敛至 $x\in\mathscr X$, 则有 $\rho(x,f(x))=d$. 
	
	如果 $d\neq 0$, 即 $x\neq f(x)$, 则 $\rho(f(x),f(f(x)))<\rho(x,f(x))=d$, 则与下确界矛盾. 所以 $d=0$.
	
	于是 $x=f(x)$, 是 $\mathscr X$ 中的不动点.
	
	如果不动点不唯一, 那么 $\rho(x_1,x_2)=\rho(f(x_1),f(x_2))<\rho(x_1,x_2)$ 矛盾.
	
	所以 $f$ 在 $\mathscr X$ 上存在唯一不动点.
	
\end{proof}

\problem[1.3.9] 设 $(M,\rho)$ 是一个紧度量空间, 又 $E\subset C(M)$, $E$ 中的函数一致有界并满足下列 Holder 条件: $$|x(t_1)-x(t_2)|\leqslant C\rho(t_1,t_2)^\alpha\quad (\forall x\in E, \forall t_1,t_2\in M),$$ 其中 $0<\alpha\leqslant 1, C>0$. 求证: $E$ 在 $C(M)$ 中是列紧集.

\begin{proof}
	根据 AA 定理, 只需证 $E$ 中函数等度连续.
	
	$\forall \varepsilon>0$, 取 $\delta=\left(\frac{\varepsilon}{C}\right)^{1/\alpha}$.
	
	那么有 $\forall \rho(t_1,t_2)<\delta, x\in E$ 有 $|x(t_1)-x(t_2)|\leqslant C\rho(t_1,t_2)^\alpha<\varepsilon$. 从而 $E$ 中函数等度连续.
	
	从而根据 AA 定理, $E$ 在 $C(M)$ 中是列紧集.
	
	
\end{proof}
